\chapter{Обыкновенные дроби}
\begin{definition}\term{Обыкновенная дробь} (\emph{common fraction}) \(\frac{a}{b}\) — запись частного \(a : b\), где \(a\) и \(b\) — целые числа и \(b \ne  0\).\end{definition}
\begin{definition}\term{Числитель} (\emph{numerator}) дроби \(\frac{a}{b}\) — число \(a\), записанное над дробной чертой.\end{definition}
\begin{definition}\term{Знаменатель} (\emph{denominator}) дроби \(\frac{a}{b}\) — ненулевое число \(b\), записанное под дробной чертой.\end{definition}
\begin{definition}\term{Правильная положительная дробь} (\emph{proper positive fraction}) — дробь с положительными числителем и знаменателем, у которой числитель меньше знаменателя.\end{definition}
\begin{definition}\term{Неправильная положительная дробь} (\emph{improper positive fraction}) — дробь с положительными числителем и знаменателем, у которой числитель не меньше знаменателя.\end{definition}
\begin{definition}\term{Смешанное число} (\emph{mixed number}) — сумма целой части и правильной дробной части; например, \(2 \frac{1}{3} = 2 + \frac{1}{3}\).\end{definition}
\section*{Основное свойство дроби}
\(\frac{a}{b} = \frac{ak}{bk}\) при \(k \ne  0\). Числитель и знаменатель можно умножать или делить на одно и то же ненулевое число.\par
\section*{Действия с дробями}
Сложение: \(\frac{a}{b} + \frac{c}{d} = \frac{ad + bc}{bd}\), где \(b,\allowbreak d \ne  0\).\par
Вычитание: \(\frac{a}{b} - \frac{c}{d} = \frac{ad - bc}{bd}\), где \(b,\allowbreak d \ne  0\).\par
Умножение: \((\frac{a}{b}) \cdot  (\frac{c}{d}) = \frac{ac}{bd}\).\par
Деление: \((\frac{a}{b}) : (\frac{c}{d}) = \frac{ad}{bc}\); дополнительно \(c \ne  0\).\par
Сравнение при \(b,\allowbreak d > 0\): сравни \(a \cdot  d\) и \(c \cdot  b\).\par
{\small\textit{Пример: \(\frac{1}{2} + \frac{1}{3} = \frac{3}{6} + \frac{2}{6} = \frac{5}{6}\). Сокращение: \(\frac{18}{24} = \frac{3}{4}\).}}\par
\Needspace{7\baselineskip}\section*{Задачи}
\begin{enumerate}[label={},leftmargin=3.2em,itemsep=10pt]
\Needspace{3\baselineskip}\item[\textbf{3.1.}] Сократи \(\frac{18}{24}\).
\Needspace{3\baselineskip}\item[\textbf{3.2.}] Преврати \(\frac{17}{5}\) в смешанное число.
\Needspace{3\baselineskip}\item[\textbf{3.3.}] Преврати \(2 \frac{3}{7}\) в неправильную дробь.
\Needspace{3\baselineskip}\item[\textbf{3.4.}] Сравни \(\frac{5}{8}\) и \(\frac{7}{12}\).
\Needspace{3\baselineskip}\item[\textbf{3.5.}] Вычисли \(\frac{3}{4} - \frac{5}{6}\).
\Needspace{3\baselineskip}\item[\textbf{3.6.}] Вычисли \(\frac{7}{15} + \frac{2}{9}\).
\Needspace{3\baselineskip}\item[\textbf{3.7.}] Вычисли \((- \frac{3}{7}) \cdot  (\frac{14}{9})\).
\Needspace{3\baselineskip}\item[\textbf{3.8.}] Вычисли \((\frac{2}{3}) : (\frac{4}{9})\).
\Needspace{3\baselineskip}\item[\textbf{3.9.}] Вычисли \(2 \frac{1}{3} - 1 \frac{5}{6}\).
\Needspace{3\baselineskip}\item[\textbf{3.10.}] Вычисли \((\frac{1}{2} + \frac{1}{3}) : (\frac{5}{6})\).
\Needspace{3\baselineskip}\item[\textbf{3.11.}] Вычисли \(\frac{\frac{3}{4}}{\frac{9}{8}}\).
\Needspace{3\baselineskip}\item[\textbf{3.12.}] Из \(80\) страниц прочитали \(\frac{3}{5}\). Сколько страниц осталось?
\Needspace{3\baselineskip}\item[\textbf{3.13.}] Из жизни. На пирог ушло \(\frac{3}{4}\) кг муки, а на блины \(\frac{2}{3}\) кг. Сколько всего муки израсходовали?
\Needspace{3\baselineskip}\item[\textbf{3.14.}] Из жизни. В бутылке \(\frac{3}{2}\) л сока. Его разливают по стаканам вместимостью \(\frac{1}{4}\) л. Сколько полных стаканов получится?
\Needspace{3\baselineskip}\item[\textbf{3.15.}] Из жизни. Турист прошёл \(\frac{2}{5}\) маршрута длиной \(15\) км. Сколько километров ему осталось?
\Needspace{3\baselineskip}\item[\textbf{3.16.}\starmark] Вычисли: \((\frac{3}{4} - \frac{5}{6}) : (\frac{7}{9} - \frac{2}{3}) + \frac{1}{2}\).
\Needspace{3\baselineskip}\item[\textbf{3.17.}\starmark] В первый день прочитали \(\frac{1}{3}\) книги, во второй — \(\frac{2}{5}\) оставшихся страниц. После этого осталось \(72\) страницы. Сколько страниц в книге?
\Needspace{3\baselineskip}\item[\textbf{3.18.}\starmark] Бак был полон. Израсходовали \(\frac{1}{4}\) воды, затем \(\frac{1}{3}\) оставшейся. После долива \(15\) л бак снова стал полным. Какова вместимость бака?
\Needspace{3\baselineskip}\item[\textbf{3.19.}] Для целых \(a,b,c,d\), где \(b\ne0\), \(d\ne0\), докажи критерий равенства дробей: \(\frac ab=\frac cd\) тогда и только тогда, когда \(ad=bc\).
\Needspace{3\baselineskip}\item[\textbf{3.20.}] Докажи, что для любого \(n\in\mathbb N\) выполнено \(\sum_{k=1}^{n}\frac1{k(k+1)}=\frac n{n+1}\).
\end{enumerate}
